\section{Introduction}

Inflammasomes are multiprotein cytosolic signaling platforms that respond to pathogen invasion, genotoxic stress, and a wide range of danger signals, and they represent a cornerstone of innate immune defense in mammals~\cite{broz2016inflammasomes,tartey2019inflammasomes}. A germline-encoded pattern recognition receptor, such as the AIM2-like receptor (ALR) family members AIM2 and IFI16 or the NOD-like receptor (NLR) family members NLRP3 and NLRP6, nucleates filamentous polymers through its pyrin domain (PYD), which in turn seed PYD filaments of the bipartite adaptor ASC; ASC then templates caspase-1 activation via CARD-CARD filaments, driving IL-1$\beta$ and IL-18 maturation and pyroptotic cell death~\cite{lu2014unified,hornung2009aim2,fernandesalnemri2009aim2,roberts2009hin200}. This sequential, prion-like assembly is extraordinarily robust: each step amplifies the incoming signal and the polymerized state persists once nucleated, yielding a binary on/off response ideally suited for rapid host defense~\cite{cai2014prion,franklin2014adaptor,matyszewski2018digital}. The same properties, however, render persistent inflammasome activity a liability. Dysregulated assembly has been implicated in cryopyrin-associated periodic syndromes and a broader spectrum of autoinflammatory disease, cancer-associated inflammation, and severe COVID-19~\cite{tartey2019inflammasomes,sharma2018pyrin,vora2021inflammasome,junqueira2022sarscov2}. Termination and tuning of inflammasome assembly are therefore biologically essential.

A major class of endogenous inflammasome regulators are the mammalian pyrin-only-proteins (POPs), small six-helix bundles of roughly 100 residues that bear a PYD but lack the downstream effector domains present in sensor or adaptor proteins~\cite{indramohan2018cops}. Three POPs (POP1, POP2, and POP3) have been characterized in humans, and they are broadly presumed to inhibit inflammasome assembly by engaging specific PYD-containing partners whose identity is inferred from primary sequence similarity. POP1 is 65\% identical to the PYD of ASC and has been proposed to block ASC nucleation directly~\cite{dealmeida2015pyrin}. POP2 is 68\% identical to the PYD of NLRP2 and has been implicated in inhibiting both ASC-dependent inflammasomes and AIM2~\cite{periasamy2017pop2,ratsimandresy2017pop2}. POP3 is 67\% identical to the PYD of AIM2 and has been reported to competitively restrict AIM2- and IFI16-driven responses to viral dsDNA~\cite{khare2014pop3}. This tripartite assignment forms the dominant textbook model of POP target specificity.

These assignments, however, rest largely on indirect inference. In cellulo and in vivo measurements of POP activity necessarily confound direct PYD-PYD recognition with the many parallel protein-protein interactions of a full inflammasome cascade, and no biochemical assay has yet resolved whether a given POP blocks nucleation, elongation, or both for a defined PYD substrate. Moreover, the premise that amino-acid sequence homology predicts POP targets is difficult to reconcile with what is now known about the architecture of PYD filaments: structural work on ASC, AIM2, NLRP3, and NLRP6 filaments has shown that every PYD engages its neighbours through six largely distinct protomer interfaces of Type 1, Type 2, and Type 3 character, and that receptor PYDs with essentially no pairwise sequence homology all converge on the single adaptor ASCPYD~\cite{lu2014unified,shen2019nlrp6,hochheiser2022directionality,matyszewski2021distinct,morrone2015assembly,sharif2019nek7,andreeva2021nlrp3}. This degenerate, code-like recognition implies that local interface complementarity, rather than global sequence identity, is what sets specificity.

We therefore set out to test the sequence-homology model of POP specificity directly, and to ask how favourable and unfavourable protomer-level interactions combine to permit or prevent PYD filament assembly. Our strategy layers three orthogonal readouts. First, Rosetta InterfaceAnalyzer computations on honeycomb-like sideviews of ASC, AIM2, IFI16, NLRP3, and NLRP6 PYD filaments quantify the interface energies ($\Delta G$, in Rosetta Energy Units) for homotypic PYD-PYD contacts and for all hypothetical POP-PYD substitutions, using the cryo-EM and crystal structures reported for each filament or its closest structural homologue~\cite{lu2014unified,shen2019nlrp6,hochheiser2022directionality,sharif2019nek7,matyszewski2021distinct}. Second, HEK293T cells, which lack endogenous inflammasome components and POPs, provide a clean co-transfection context in which mCherry-tagged PYDs and eGFP-tagged POPs can be scored for filament or punctum suppression in isolation. Third, FRET-based polymerization kinetics with purified recombinant POPs and PYDs, complemented by negative-stain electron microscopy (nsEM), dissect whether a given POP blocks the rate-limiting nucleation step, attenuates elongation, or both~\cite{mazanek2019fret,matyszewski2018digital,morrone2015assembly}.

The central contributions of this work are: (i) we show that POP1 is a poor inhibitor of ASCPYD ($\leq 20\%$ filament suppression in HEK293T cells at 1200~ng POP-eGFP and no detectable effect on in vitro polymerization up to 150~$\mu$M), overturning the canonical POP1-as-ASC-inhibitor model~\cite{dealmeida2015pyrin}; (ii) we demonstrate that POP1 instead acts predominantly on upstream receptor PYDs (AIM2, IFI16, NLRP3, and NLRP6), suppressing their filament assembly by 60--70\% in cells; (iii) we show that POP2 and POP3 potently block ASCPYD nucleation, extending the lag phase by up to $\sim$15~min in FRET experiments and abolishing or strongly reducing ASCPYD filaments in nsEM; and (iv) we put forward a design principle in which inhibitory POPs must simultaneously present at least one favourable recognition interface and at least one strongly unfavourable (repulsive) interface with their PYD substrate. Together these results redefine POPs as decoy adaptors and receptors, demote sequence homology as a predictor of function, and provide a protomer-interface framework that can guide the rational design of synthetic POP variants or small-molecule mimics for therapeutic targeting of persistent inflammasome activity.

\section{Related Work and Background}

\paragraph{Inflammasome architecture and prion-like assembly.} Early biochemical and cryo-EM work established that ASC-dependent inflammasomes assemble through successive PYD and CARD filaments in a unified template mechanism: activated receptors nucleate ASCPYD filaments, which in turn bundle CARD filaments of caspase-1~\cite{lu2014unified,shen2019nlrp6,sharif2019nek7,andreeva2021nlrp3,hochheiser2022directionality,hochheiser2022seeding}. The stability of these filaments, the absence of a defined off-state, and their capacity to convert soluble subunits across cells motivated the recognition of ASC as a functional prionoid adaptor~\cite{cai2014prion,franklin2014adaptor}. Quantitative kinetic dissection using FRET-based polymerization assays has since shown that assembly proceeds through a slow nucleation phase followed by fast elongation, that longer dsDNA stimuli accelerate AIM2 nucleation, and that the resulting digital circuit enforces ultrasensitive signaling~\cite{matyszewski2018digital,mazanek2019fret,morrone2015assembly}. Our work extends this framework by asking how a small inhibitor protein can perturb either the nucleation or the elongation phase without co-assembling into the mature filament.

\paragraph{ALR and NLR receptor PYDs.} AIM2 and IFI16 are the central cytosolic dsDNA sensors of the ALR family, with AIM2 forming inflammasome complexes in response to pathogen and self dsDNA and IFI16 contributing to sensing of nuclear-replicating viruses such as KSHV~\cite{hornung2009aim2,fernandesalnemri2009aim2,roberts2009hin200,kerur2011ifi16,morrone2015assembly}. NLRP3 and NLRP6 contain tripartite NACHT-LRR-PYD architectures and have been implicated in sterile and microbial inflammation, respectively~\cite{shen2019nlrp6,sharif2019nek7,andreeva2021nlrp3}. The PYDs of these receptors share a common six-helix fold but, unlike the adaptor ASCPYD, show limited pairwise sequence identity among themselves, yet all funnel into ASCPYD polymerization--an empirical observation that motivates the degenerate-code-like recognition model we invoke~\cite{matyszewski2021distinct,lu2014unified}. Both AIM2 and IFI16 have also been identified as autoantigens in systemic lupus erythematosus, underscoring the clinical relevance of tight regulation of ALR-dependent signaling~\cite{antiochos2022dna}.

\paragraph{Endogenous inhibitors: POPs and COPs.} The mammalian repertoire of endogenous inflammasome inhibitors includes PYD-only proteins (POP1, POP2, POP3) and CARD-only proteins (COPs), all of which interfere with death-domain oligomerization~\cite{indramohan2018cops}. Based primarily on in cellulo co-transfection assays and transgenic mouse models, POP1 has been assigned ASC as its principal target, POP2 has been linked to both ASC and NLRP inflammasomes including a proposed primary target of NLRP2PYD, and POP3 has been assigned to ALRs~\cite{dealmeida2015pyrin,periasamy2017pop2,ratsimandresy2017pop2,khare2014pop3}. These assignments, however, leave the underlying biophysical mechanism unspecified. In particular, none of the prior studies distinguish whether inhibition occurs at the nucleation step, at the elongation step, or through stoichiometric competition with endogenous PYDs; and none directly measure the interface energetics that govern PYD-PYD versus POP-PYD recognition. Our combined Rosetta, cellular, and biochemical analysis fills this gap and, as we show below, revises the canonical POP--target pairings.

\paragraph{Clinical relevance.} Dysregulated inflammasome activity has been implicated in a growing list of human pathologies including autoinflammatory syndromes, inflammatory bowel disease, and cancer-associated inflammation~\cite{tartey2019inflammasomes,sharma2018pyrin}. Severe COVID-19 has in particular been linked to NLRP3- and AIM2-driven pyroptotic inflammation in infected monocytes~\cite{vora2021inflammasome,junqueira2022sarscov2}. A mechanistic understanding of how endogenous POPs throttle these pathways is therefore not only biochemically informative but also therapeutically actionable, as it supplies a structural template for engineered or pharmacological mimics.
